\chapter{Fasciotomy}

\section*{What Is This Surgery?}
Fasciotomy is a critical surgical procedure involving the incision and opening of the inelastic fascial compartments within a limb. This intervention is performed to relieve pathologically elevated tissue pressure that, if left untreated, would compromise the viability of muscles, nerves, and blood vessels, leading to irreversible ischemic damage. Primarily indicated for acute compartment syndrome, a time-sensitive orthopedic emergency, fasciotomy aims to restore adequate blood flow and prevent permanent functional impairment or limb loss. The surgery is most commonly performed on the lower leg and forearm, but can be applied to any anatomical region encased by restrictive fascia, such as the thigh, hand, or foot, where compartment syndrome may develop. Its clinical significance lies in its ability to rapidly decompress compromised tissues, thereby salvaging limb function and preventing devastating long-term complications.

\section*{Alternative Names}
\begin{itemize}
  \item Compartment release
  \item Decompressive fasciotomy
  \item Compartment syndrome surgery
  \item Limb decompression
  \item Four-compartment fasciotomy (specifically for the lower leg)
  \item Volar fasciotomy (specifically for the forearm)
\end{itemize}

\section*{Relevant Surgical Anatomy}
The lower leg, a common site for compartment syndrome, is encased by a dense, inelastic deep fascia and divided into four distinct compartments: anterior, lateral, superficial posterior, and deep posterior. Each compartment contains specific muscle groups, nerves, and blood vessels, making a thorough understanding of their boundaries and contents crucial for complete surgical decompression. The anterior compartment houses the tibialis anterior, extensor digitorum longus, extensor hallucis longus, and fibularis tertius muscles, along with the deep peroneal nerve and anterior tibial artery and vein. The lateral compartment contains the fibularis longus and brevis muscles, innervated by the superficial peroneal nerve.

The superficial posterior compartment comprises the gastrocnemius, soleus, and plantaris muscles, drained by the small saphenous vein and supplied by branches of the tibial nerve. Deep to the soleus, the deep posterior compartment contains the tibialis posterior, flexor digitorum longus, and flexor hallucis longus muscles, along with the posterior tibial artery and vein, and the tibial nerve. The peroneal artery and vein also lie within this compartment. Surgical landmarks for the lower leg include the palpable tibial crest anteriorly, the fibula laterally, and the medial border of the tibia posteromedially. In the forearm, the major compartments include the volar (flexor) and dorsal (extensor) compartments, with critical neurovascular structures such as the median, ulnar, and radial nerves, and their corresponding arteries, requiring meticulous protection during fascial release.

\section{Surgical Principle \& Rationale}
The fundamental principle behind fasciotomy is the immediate reduction of elevated intracompartmental pressure to restore tissue perfusion. In compartment syndrome, an increase in fluid volume (e.g., from hemorrhage, edema, or reperfusion injury) within a closed, unyielding fascial compartment leads to a rise in interstitial pressure. When this pressure exceeds the capillary perfusion pressure, blood flow to the muscles and nerves within that compartment is compromised, leading to ischemia. Prolonged ischemia results in cellular anoxia, muscle necrosis (rhabdomyolysis), and irreversible nerve damage.

The surgical division of the restrictive fascia effectively increases the volume of the compartment, thereby decreasing the intracompartmental pressure. This decompression allows for the re-establishment of arterial inflow and venous outflow, restoring oxygen and nutrient delivery to the threatened tissues. The technical goals of the procedure are to achieve a complete and thorough release of all affected fascial compartments, ensuring that no constricting bands remain, and to preserve the overlying skin and subcutaneous tissues as much as possible for eventual wound closure.

\section{History \& Surgical Evolution}
The concept of compartment syndrome and the need for surgical decompression has evolved over centuries. Early descriptions of ischemic contractures, now known as Volkmann's ischemic contracture, date back to the 19th century, notably by Richard von Volkmann in 1881, who attributed the condition to tight bandages and arterial occlusion. However, the understanding of elevated intracompartmental pressure as the primary etiology and the systematic application of fasciotomy as a treatment gained prominence in the 20th century.

During World War I and II, the high incidence of vascular injuries and subsequent reperfusion edema led to a greater appreciation for the role of fasciotomy in preventing limb loss and preserving function. Surgeons like Seddon and Ellis further refined the understanding of the pathophysiology and the surgical techniques. The development of specific approaches, such as the two-incision four-compartment fasciotomy for the lower leg, became standard practice in the latter half of the 20th century, significantly improving outcomes for patients with acute compartment syndrome. The advent of objective compartment pressure measurement devices further aided in the timely diagnosis and indication for this critical procedure.

\section{Clinical Indications}
Fasciotomy is primarily indicated for conditions characterized by dangerously elevated intracompartmental pressure that threatens limb viability. Key clinical indications include:
\begin{itemize}
  \item Acute compartment syndrome resulting from traumatic injuries such as long bone fractures (especially tibial shaft fractures), crush injuries, or severe soft tissue trauma.
  \item Reperfusion injury following prolonged limb ischemia, such as after arterial repair for vascular trauma or embolectomy.
  \item Severe burns, particularly circumferential full-thickness burns, which can create an unyielding eschar that restricts tissue expansion and elevates pressure.
  \item High-pressure injection injuries, where injected substances can rapidly increase compartment volume and pressure.
  \item Prolonged limb compression, often seen in unconscious patients (e.g., drug overdose, prolonged surgery in certain positions) leading to muscle necrosis and swelling.
  \item Tight casts, splints, or dressings that inadvertently constrict a limb and prevent normal tissue expansion.
  \item Necrotizing soft tissue infections (e.g., necrotizing fasciitis), where fasciotomy is performed not only for decompression but also for extensive debridement of infected and necrotic tissue.
  \item Chronic exertional compartment syndrome, typically in athletes, where repetitive muscle activity causes transient but symptomatic pressure elevation, usually managed electively after conservative measures fail.
\end{itemize}

\section*{Clinical Positioning}
For acute compartment syndrome, fasciotomy is considered the primary and definitive emergency treatment. Once the diagnosis is made, either clinically or confirmed by compartment pressure measurements, the procedure should not be delayed, as time is of the essence to prevent irreversible muscle and nerve damage. In this context, it is a limb-salvaging intervention, serving as the first-line surgical approach to restore perfusion and prevent permanent disability.

For necrotizing soft tissue infections, fasciotomy serves as a primary surgical intervention for source control and debridement, often performed urgently alongside aggressive antibiotic therapy. While it is a critical initial step, it may be followed by serial debridements and reconstructive procedures, making it part of a multi-stage treatment plan. In cases of chronic exertional compartment syndrome, fasciotomy is an elective procedure, typically performed after conservative measures have failed, and is considered a primary surgical solution for symptom relief and return to activity.

\section{Surgical Technique \& Procedure}
Fasciotomy is typically performed under general anesthesia, though regional anesthesia may be used in select cases. The patient is positioned to allow full access to the affected limb, which is then prepped and draped in a sterile fashion. The specific surgical approach varies depending on the limb involved, with the lower leg and forearm being the most common sites.

For the lower leg, the standard approach is the two-incision four-compartment fasciotomy:
\begin{enumerate}
  \item \textbf{Anesthesia and Positioning:} General anesthesia is induced. The patient is typically placed in a supine position, and the affected leg is prepped from the foot to the thigh, allowing for full visualization and access.
  \item \textbf{Anterolateral Incision:} A longitudinal skin incision, approximately 15-20 cm in length, is made on the anterolateral aspect of the lower leg, midway between the tibial crest and the fibula. This incision extends through the skin and subcutaneous tissue.
  \item \textbf{Anterior Compartment Release:} The fascia overlying the anterior compartment is identified and incised longitudinally for its entire length, from just distal to the knee joint to just proximal to the ankle joint. Care is taken to protect the deep peroneal nerve, which lies within this compartment.
  \item \textbf{Lateral Compartment Release:} The fascia overlying the lateral compartment is then identified, typically posterior to the anterior compartment release, and incised longitudinally for its full length. The superficial peroneal nerve, which traverses this compartment, must be carefully identified and protected.
  \item \textbf{Posteromedial Incision:} A second longitudinal skin incision, similar in length, is made on the posteromedial aspect of the lower leg, approximately 2 cm posterior to the palpable medial border of the tibia. This incision also extends through the skin and subcutaneous tissue.
  \item \textbf{Superficial Posterior Compartment Release:} The fascia of the superficial posterior compartment is identified and incised longitudinally. Care is taken to protect the saphenous vein and nerve, which are superficial in this region.
  \item \textbf{Deep Posterior Compartment Release:} The soleus muscle is then retracted posteriorly to expose the deep posterior compartment fascia. This fascia is also incised longitudinally for its entire length. Extreme caution is exercised to protect the posterior tibial neurovascular bundle (posterior tibial artery, vein, and tibial nerve), which lies within this compartment.
  \item \textbf{Assessment and Hemostasis:} After all four compartments are completely released, the muscles are inspected for viability (e.g., color, contractility, bleeding). Any obviously necrotic tissue is debrided. Meticulous hemostasis is achieved.
  \item \textbf{Wound Management:} The skin incisions are typically left open due to anticipated postoperative swelling. The wounds are loosely packed with sterile, non-adherent dressings, and the limb is elevated.
\end{enumerate}
For the forearm, a volar incision (often Henry's approach) is used to release the flexor compartments, and a dorsal incision is used for the extensor compartments. Protection of the median, ulnar, and radial nerves and vessels is paramount. The wounds are typically left open and managed with serial dressing changes. Delayed primary closure, often with split-thickness skin grafting, is performed days to weeks later once swelling has subsided and the wounds are clean and granulating. Drains are generally not used in fasciotomy wounds, as they are left open.

\section*{Preoperative Workup \& Preparation}
Given the emergent nature of acute compartment syndrome, preparation for fasciotomy is often rapid and focused on immediate stabilization and diagnosis. Clinical suspicion is paramount, and if time permits and the patient is cooperative, compartment pressure measurements may be performed to objectively confirm the diagnosis, though clinical signs often dictate immediate surgery. A rapid medical and cardiac evaluation is performed to assess the patient's fitness for general anesthesia.

Key preoperative steps include:
\begin{itemize}
  \item \textbf{Laboratory Tests:} Routine preoperative blood work includes a complete blood count (CBC), basic metabolic panel (BMP) to assess renal function and electrolytes (especially potassium and creatinine kinase due to potential rhabdomyolysis), coagulation studies (PT/INR, PTT), and type and cross-match for potential blood transfusion.
  \item \textbf{Imaging:} X-rays of the affected limb are obtained to identify any associated fractures. Doppler ultrasound or CT angiography may be used to assess vascular integrity if a vascular injury is suspected.
  \item \textbf{NPO Status:} The patient is kept NPO (nil per os) to minimize aspiration risk during anesthesia.
  \item \textbf{Medication Review:} Anticoagulants or antiplatelet medications may need to be reversed or held, though the urgency of fasciotomy often outweighs this concern.
  \item \textbf{Prophylactic Antibiotics:} Intravenous broad-spectrum antibiotics are administered preoperatively to reduce the risk of surgical site infection, especially given the open wounds.
  \item \textbf{Informed Consent:} Despite the emergency, informed consent is obtained, explaining the procedure, its necessity for limb salvage, and potential risks including nerve damage, infection, and the need for further surgeries or skin grafting.
  \item \textbf{Anesthesia Consultation:} Anesthesia is consulted immediately for rapid assessment and preparation for general or regional anesthesia.
\end{itemize}

\section*{Contraindications \& Precautions}
\textbf{Absolute Contraindications:}
\begin{itemize}
  \item \textbf{Irreversible Muscle Necrosis:} If compartment syndrome has been present for an extended period (typically >6-8 hours) and irreversible muscle necrosis has already occurred, fasciotomy may not be beneficial and could increase the risk of infection and systemic complications (e.g., sepsis from necrotic tissue, acute kidney injury from rhabdomyolysis). In such cases, primary amputation might be considered.
\end{itemize}

\textbf{Relative Contraindications and Precautions:}
\begin{itemize}
  \item \textbf{Severe Coagulopathy:} While not an absolute contraindication in an emergency, severe uncorrected coagulopathy increases the risk of significant bleeding. Efforts should be made to correct it preoperatively if time allows, but this should not delay an urgent fasciotomy.
  \item \textbf{Severe Systemic Illness:} Patients with severe comorbidities may have increased anesthetic and surgical risks. The decision to proceed must weigh the benefits of limb salvage against the patient's overall medical stability.
  \item \textbf{Non-Viable Limb:} If the limb is clearly non-viable due to extensive trauma or prolonged ischemia beyond the point of salvage, primary amputation may be a more appropriate and safer course of action.
  \item \textbf{Delay in Treatment:} The most critical precaution is to avoid any delay once acute compartment syndrome is diagnosed. Every hour of delay significantly increases the risk of permanent muscle and nerve damage.
  \item \textbf{Incomplete Release:} Care must be taken to ensure all affected compartments are completely released along their entire length to prevent persistent pressure and ongoing ischemia.
  \item \textbf{Neurovascular Protection:} Surgeons must exercise extreme caution to identify and protect vital neurovascular structures during fascial incisions, particularly the superficial and deep peroneal nerves in the lower leg, and the median, ulnar, and radial nerves in the forearm.
\end{itemize}

\section*{Complications \& Risks}
Fasciotomy, while limb-saving, carries several potential risks and complications:
\begin{itemize}
  \item \textbf{General Surgical Risks:}
    \begin{itemize}
      \item \textbf{Bleeding:} Hematoma formation, significant blood loss requiring transfusion.
      \item \textbf{Infection:} Surgical site infection, cellulitis, osteomyelitis, sepsis.
      \item \textbf{Anesthesia Complications:} Adverse reactions to anesthetic agents, respiratory or cardiac complications, aspiration.
    \end{itemize}
  \item \textbf{Procedure-Specific Risks:}
    \begin{itemize}
      \item \textbf{Nerve Injury:} Damage to superficial or deep nerves (e.g., superficial peroneal nerve, saphenous nerve, posterior tibial nerve) leading to sensory deficits (numbness, paresthesia) or motor weakness/paralysis (e.g., foot drop).
      \item \textbf{Muscle Damage:} Direct injury to muscle tissue during incision, or persistent necrosis if decompression is delayed or incomplete.
      \item \textbf{Cosmetic Scarring:} Significant, often wide, and disfiguring scars due to the large incisions and delayed closure, which may require skin grafting.
      \item \textbf{Chronic Pain:} Persistent pain at the incision sites or due to nerve irritation or neuroma formation.
      \item \textbf{Wound Complications:} Wound dehiscence, delayed wound healing, need for extensive skin grafting or flap coverage, hypertrophic scarring.
      \item \textbf{Contracture:} Volkmann's ischemic contracture, characterized by muscle fibrosis and shortening, leading to permanent deformity and functional loss, especially if decompression is delayed.
      \item \textbf{Rhabdomyolysis and Renal Failure:} Release of muscle breakdown products (myoglobin) into the bloodstream can lead to acute kidney injury, particularly if significant muscle necrosis occurred prior to decompression.
      \item \textbf{Limb Loss:} Despite fasciotomy, severe, prolonged ischemia or overwhelming infection can still necessitate amputation.
      \item \textbf{Deep Vein Thrombosis (DVT) / Pulmonary Embolism (PE):} Risk associated with immobility and surgery, requiring prophylactic measures.
      \item \textbf{Recurrence of Compartment Syndrome:} Rare, but possible if the initial release was incomplete or if there is re-bleeding into the compartment.
    \end{itemize}
\end{itemize}

\section*{Postoperative Recovery Timeline}
Immediately after fasciotomy, the patient is transferred to the Post-Anesthesia Care Unit (PACU) for close monitoring of vital signs, pain levels, and neurovascular status of the affected limb. Pain management is a priority, often involving intravenous analgesics, and the limb is typically elevated to reduce swelling. Neurovascular checks (assessing sensation, motor function, capillary refill, and pulses) are performed frequently to ensure adequate perfusion and detect any new deficits. The open fasciotomy wounds are covered with sterile, non-adherent dressings, often loosely packed to allow for continued drainage and swelling. The first 24-48 hours are critical for monitoring for signs of re-bleeding, infection, or persistent compartment syndrome.

Over the next several days, the patient will undergo serial wound examinations and dressing changes. Depending on the extent of muscle damage and the presence of infection, repeat surgical debridements may be necessary to remove any non-viable tissue. Once the swelling has subsided and the wounds are clean and granulating, typically within 3 to 7 days, definitive wound closure is performed. This may involve delayed primary closure (suturing the skin edges together if there is minimal tension), or more commonly, split-thickness skin grafting to cover the wide fascial defects. Physical therapy and rehabilitation are initiated early to prevent joint stiffness and muscle atrophy, gradually progressing as the wounds heal over several weeks to months. Full recovery, including return to strenuous activities, can take 3-6 months or longer, depending on the severity of the initial injury and any residual deficits.

\section*{Surgical Findings \& Pathology}
During fasciotomy, the surgeon makes several key observations that guide the procedure and inform the prognosis. Upon incision of the skin and subcutaneous tissue, the underlying fascia is typically found to be tense and non-yielding. Once the fascia is incised, the muscles within the compartment may bulge significantly, indicating the release of high pressure. The surgeon assesses muscle viability by observing its color (healthy muscle is typically reddish-pink, while ischemic muscle may appear pale, dusky, or dark red), contractility (response to stimulation), and bleeding upon incision. The presence of hematoma, edema, or obviously necrotic tissue is noted and debrided as necessary. Confirmation of complete fascial release in all affected compartments is paramount.

While not routinely performed, biopsies of muscle tissue may be taken if viability is uncertain or if there is suspicion of other underlying pathologies. Pathology reports on resected tissues would typically describe findings consistent with muscle fiber necrosis, interstitial edema, hemorrhage, and inflammatory infiltrate, confirming the diagnosis of ischemic injury. In cases of necrotizing fasciitis, the pathology report would detail extensive tissue necrosis, bacterial presence, and inflammatory changes, guiding further debridement. These findings, combined with clinical assessment, help determine the extent of damage and guide postoperative management.

\section{Expected Postoperative Course}
A normal and successful postoperative course following fasciotomy signifies the effective resolution of compartment syndrome and the prevention of permanent ischemic injury. Patients typically experience a significant reduction in the severe pain associated with compartment syndrome shortly after decompression. Over the first few days, swelling in the limb should gradually subside, allowing for eventual wound closure. Neurovascular function, including sensation and motor strength, should progressively improve, although some residual deficits may persist depending on the duration and severity of ischemia prior to surgery. The fasciotomy wounds will require meticulous care, with serial dressing changes and eventual delayed closure or skin grafting. Patients will gradually regain mobility and strength through a structured physical therapy program. A successful course is characterized by the absence of major complications such as infection, significant nerve damage, or the development of contractures, leading to maximal functional recovery and return to daily activities.

\section{Postoperative Care \& Follow-up}
Post-fasciotomy care involves a structured follow-up plan to ensure optimal recovery and address any potential complications. This comprehensive approach is crucial for maximizing functional outcomes.

Key aspects of postoperative care and follow-up include:
\begin{itemize}
  \item \textbf{Serial Wound Care:} Regular clinic visits for wound checks, dressing changes, and monitoring for signs of infection or delayed healing. This continues until definitive wound closure is achieved.
  \item \textbf{Suture/Staple Removal:} If delayed primary closure was performed, sutures or staples are typically removed 2-3 weeks post-closure.
  \item \textbf{Physical and Occupational Therapy:} Crucial for regaining strength, range of motion, and function. Therapy often begins early in the post-operative period, sometimes even before definitive wound closure, and continues for several months.
  \item \textbf{Scar Management:} Techniques such as massage, silicone sheeting, or compression garments may be recommended to improve scar appearance and reduce contracture formation.
  \item \textbf{Activity Resumption:} Gradual return to normal activities, guided by the surgeon and physical therapist, avoiding strenuous activities until adequate healing and strength are achieved.
  \item \textbf{Warning Signs:} Patients are educated to report any signs of infection (increased redness, swelling, pus, fever), worsening pain, new numbness or weakness, or signs of deep vein thrombosis (calf pain, swelling).
  \item \textbf{Reconstructive Procedures:} In cases of significant tissue loss, nerve damage, or severe contracture, later reconstructive surgeries (e.g., tendon transfers, nerve grafts, flap coverage) may be necessary to improve function and cosmesis.
\end{itemize}
Long-term follow-up ensures ongoing assessment of functional recovery and addresses any persistent issues.

\section*{Epidemiology}
Acute compartment syndrome (ACS) is a relatively uncommon but critical orthopedic emergency, with an estimated incidence ranging from 3 to 30 per 100,000 population per year, varying significantly based on the prevalence of trauma. It predominantly affects young males, often associated with high-energy trauma such as motor vehicle accidents or sports injuries. Tibial shaft fractures are the most common cause of ACS in the lower leg, accounting for approximately 30-75\% of cases. Other significant causes include crush injuries, reperfusion injury after vascular repair, and severe burns. Chronic exertional compartment syndrome (CECS) is more prevalent in athletes, particularly runners and military personnel, with an incidence that is harder to quantify but is a common cause of exercise-induced leg pain. While the mortality rate for isolated ACS is low, the morbidity associated with delayed diagnosis and treatment is substantial, leading to permanent functional deficits in up to 50\% of cases, including contracture, sensory loss, and muscle weakness, and in severe cases, limb amputation.

\section*{Outcomes \& Prognosis}
The outcome and prognosis following fasciotomy are highly dependent on the timeliness of diagnosis and intervention. When performed promptly (ideally within 6 hours of symptom onset), fasciotomy has a high success rate in preventing irreversible ischemic damage and preserving limb viability. Functional outcomes are generally excellent in these cases, with many patients achieving near-normal function, although some may experience residual scarring, mild sensory deficits, or muscle weakness. Delayed decompression, however, significantly worsens the prognosis, leading to a higher incidence of permanent sequelae such as Volkmann's ischemic contracture, foot drop, chronic pain, and the need for further reconstructive surgeries. Recurrence of acute compartment syndrome after a complete fasciotomy is rare. For chronic exertional compartment syndrome, fasciotomy typically provides good to excellent relief of symptoms in 70-90\% of patients, allowing them to return to their desired activity levels. Prognostic factors include the severity of the initial injury, the duration of ischemia before decompression, the completeness of the fascial release, and the patient's adherence to postoperative rehabilitation.

\section*{References}
\begin{enumerate}
  \item MedlinePlus. Fasciotomy. U.S. National Library of Medicine; 2024. Available from: \texttt{https://medlineplus.gov/ency/article/007222.htm}
  \item Schwartz's Principles of Surgery. 11th ed. Brunicardi FC, Andersen DK, Billiar TR, et al., editors. McGraw-Hill Education; 2019. Chapter 48: Orthopedic Surgery.
  \item Sabiston Textbook of Surgery: The Biological Basis of Modern Surgical Practice. 21st ed. Townsend CM Jr, Beauchamp RD, Evers BM, Mattox KL, editors. Elsevier; 2021. Chapter 48: Orthopedic Surgery.
  \item American Academy of Orthopaedic Surgeons (AAOS). Clinical Practice Guideline: Management of Acute Compartment Syndrome. 2020.
\end{enumerate}

\clearpage